\section{原创贡献}

本论文聚焦\gls{rl}场景中特有的两类延迟：状态观测延迟（state observation delay）和动作执行延迟（action execution delay）。目标是拓宽对此类延迟的理解并提供高效解决方案。特别关注文献中普遍存在且在许多情况下能很好近似实际延迟的常值延迟（constant delay）。然而，我们将反复研究将框架扩展到更原始延迟类型的理论和实证影响。本论文的贡献遵循四个目标：提供\gls{rl}中延迟的统一框架；为更好理解延迟提供理论基础；设计有理论依据且高效的算法以应对延迟环境；对所提算法及文献中算法进行广泛实证评估。下文将详述每个目标的贡献，然后引用本人合著的相关出版物。

\subsubsection{延迟的统一框架}

本文对\gls{rl}及相关领域的延迟文献进行了广泛综述。重点放在\gls{rl}中的状态观测和动作执行延迟，但也讨论了其他类型延迟。值得注意的是，我们将文献中的算法归纳为三类。第一类是\emph{增广状态（augmented state）}方法。其思想是将最后观测到的状态与智能体已选择但尚未观察到结果的动作序列进行增广。使用这种新状态可将延迟问题重新转化为传统的无延迟\gls{mdp}框架。第二类是\emph{无记忆（memoryless）}方法，智能体忽略延迟，仅使用最后观测状态的信息。这一思想简单，但在实践中可能高效。最后，\emph{基于模型（model-based）}方法学习环境模型以预测动作将施加时的环境状态。据我们所知，此前尚无工作提供\gls{rl}中延迟的综述。我们认为提供此信息对该领域有益，因为文献中经常出现重叠结果。

此外，我们提出延迟问题的统一数学框架。这一提议遵循同一思路，即延迟文献可受益于更多结构化。我们的框架将所有最常見的延迟类型作为特例包含在内，并与\gls{mdp}模型自然衔接。

\subsubsection{延迟的理论理解}

首先，分析延迟状态观测或动作执行对\gls{mdp}的理论影响。第一个结果是正式证明一个看似显而易见但尚未被证明的事实：延迟越长，最大可达累积奖励（即回报）越低。同理，展示延迟如何影响策略集所能实现的回报变异性。更长的延迟缩小了智能体可获得回报的范围，但也降低了重复学习同一任务时回报的方差。然后，将常值延迟框架扩展到多动作延迟（multiple action delay）框架，其中单个动作可影响未来的多个转移。该框架从理论角度具有研究价值，因为它将常延迟过程作为特例包含在内。此外，它自然地将高阶\glspl{mc}模型扩展到\glspl{mdp}。我们展示多动作延迟过程可通过状态增广重新转化为\gls{mdp}，与常延迟过程类似。然后研究该框架相对于传统\gls{mdp}的特性和特殊性。特别是，这有助于强调高阶\gls{mdp}与延迟\gls{mdp}之间的联系。值得注意的是，一些性质（如单链性（unichain））不会从底层\gls{mdp}传递到基于它构建的延迟过程，而其他性质（如连通性（communicating））则会传递。

其次，将焦点转向延迟\gls{rl}算法并研究其保证。分析处理延迟的三种方法：增广状态、无记忆和基于模型的方法。已知增广状态方法可为延迟问题找到最优策略。相反，我们展示基于模型方法可能过于严格而丢弃最优策略。然而，在平滑\gls{mdp}（smooth MDP）——将在本文中定义的概念——背景下，我们证明基于模型方法能提供接近最优的策略。将上述保证扩展到非整数延迟（non-integer delay）。非整数延迟让人联想到实时\gls{rl}的异步环境假设。智能体以给定频率观测状态，但只能以相位变化的方式执行动作。最后，推导出基于模型策略在常值延迟上训练但在随机延迟上测试时的性能界。

\subsubsection{延迟强化学习算法}

在算法方面，首先探索基于模型方法，设计一种能够学习当前状态概率分布（即信念（belief））的向量化表示的算法。为学习信念表示，第一个神经网络（Transformer）处理增广状态并为每个未来状态生成表示（直至无延迟状态）。该表示通过拟合分布的另一个网络（归一化流（normalising flow）网络）训练以对应信念。将该表示接入\gls{rl}算法可提升性能并加速学习速度，效果在随机环境中尤为明显。此外，提供一种直接方法，使预训练于特定延迟的该算法能适应更短延迟。

然后探索一种不同的方法，通过模仿无延迟专家行为来学习策略。实际中，延迟智能体被训练为复制无延迟智能体基于当前状态选择的动作，但仅使用延迟信息——即增广状态。展示如何利用平滑\glspl{mdp}的理论保证与模仿学习（imitation learning）的理论结果相结合来指导算法设计。实证表明该方法达到最先进性能，且足够通用以适用于常值、非整数和随机延迟。

作为第三种方法，探索无记忆方法。在此情形下，问题本质上变成非平稳（non-stationary）。因此，设计一种理论，通过优化未来性能估计的概率下界来适应未来非平稳性。该未来性能估计基于多重重要性采样（importance sampling）以考虑非平稳性导致的过去动态多样性。在平滑非平稳条件下，该估计器的偏差有界。为进一步考虑过程的随机性，利用浓度不等式获得估计的下界。这涉及前一估计器的方差上界。有趣的是，控制估计器的方差间接约束了策略的非平稳性。这防止对过去非平稳动态的过拟合并更好地泛化到未来动态。实际中，算法实现为一个超策略（hyper-policy），接受时间作为输入并输出每个时刻要查询的策略。实验展示该算法如何学习过程的非平稳性以适应它，包括当非平稳性由延迟引起时。

\subsubsection{广泛实证评估}

最后一个贡献是对本文提出的算法及文献中算法进行广泛实证评估。设计多种测试任务，从平衡摆到货币汇率交易。在这些任务中，改变延迟的长度和性质以研究算法的鲁棒性。考虑常值、非整数、随机和多动作延迟。这些实证评估有助于理解算法的特性和偏好设置。

\subsubsection{相关出版物}

上述贡献可在以下论文中找到。\cite{liotet2021learning} 中提出基于模型的方法并分析其可能过于严格的问题。在该工作中，我设计了方法，进行了理论分析和大部分实验分析。模仿学习方法及平滑\glspl{mdp}中延迟的分析见 \cite{liotet2022delayed}。我完成了常值延迟的部分理论分析以及向非整数和随机延迟的扩展，并进行了大部分实验分析。非平稳无记忆方法来自 \cite{liotet2022lifelong}。在该工作中，我进行了大部分理论分析、部分方法设计以及全部实证研究。最后，关于多动作延迟扩展及改变延迟分布对智能体性能影响的分析结果将作为待提交论文的基础。在该工作中，我完成了全部理论和实证分析。
